\chapter{Recurso hídrico, demanda y energía}
\section{Estudio hidrológico necesario}
El borrador reporta área de cuenca, precipitación, caudal medio y Q95 sin estación, periodo, serie ni metodología de transferencia. La cercanía de Phara o Sandia no demuestra equivalencia hidrológica. Se requieren delimitación de cuenca en el punto de captación, registros SENAMHI/ANA con metadatos, control de faltantes, homogeneidad, aforos locales y método de regionalización calibrado.

Cada aforo debe registrar fecha, coordenadas, sección, nivel, equipo, calibración, procedimiento y error estimado. Una campaña en estiaje ayuda a contrastar el modelo, pero un aforo aislado no genera una curva de duración ni garantiza caudal todo el año. Deben revisarse derechos y captaciones aguas arriba y en el tramo derivado, retornos, avenida de diseño y transporte de sedimentos.

\section{Balance disponible y caudal de diseño}
Para cada intervalo $i$:
\begin{align}
Q_{disp,i}&=\max[0,Q_{rio,i}-Q_{eco,i}-Q_{terceros,i}-Q_{otros,i}],\\
Q_{turb,i}&=\min[Q_d,Q_{disp,i}].
\end{align}
El balance debe ser consistente con la sección de aforo: si el caudal observado ya refleja una extracción aguas arriba, no se descuenta nuevamente. Las purgas o pérdidas se incorporan donde realmente consumen o desvían agua. Cuando $Q_{turb,i}$ quede bajo el mínimo técnico del grupo, se analiza parada o funcionamiento parcial según fabricante.

Con los valores no comprobados del borrador, $0.520-0.100=0.420$ m$^3$/s. Este número coincide exactamente con el caudal de la alternativa y deja \textbf{margen nulo} para otros usos, incertidumbre y purgas. Por ello no respalda 150 kW firmes. Ni siquiera un Q95 validado representa 100\% del tiempo.

La curva de duración se obtendrá ordenando caudales de una serie definida, documentando periodo y posición de trazado (por ejemplo, $p=m/(n+1)$). Se debe diferenciar curva anual, mensual y de época seca. En esta revisión se genera una \emph{curva de potencia versus caudal}, no una falsa curva de duración del río.

\section{Demanda que justifica la central}
Hay dos problemas distintos. En paralelo a red, el objetivo puede ser exportar energía conforme a conexión y contrato; no implica atender un número inventado de viviendas. En modo aislado, se necesita inventario por usuario, demanda simultánea, perfil horario, cargas prioritarias, motores y crecimiento sustentado. Los 150 kW del borrador son una meta de oferta, no una demanda demostrada.

El levantamiento de cargas externas deberá contener usuario o grupo, cantidad, potencia eléctrica de entrada, factor de potencia, horas de uso, simultaneidad medida o justificada, arranque y prioridad. Para cada intervalo:
\[
P_{dem}(t)=\sum_j n_jP_{u,j}f_j(t),\qquad
Q_{dem}(t)=\sum_j P_j(t)\tan[\arccos(fp_j)].
\]
Se suman potencias reactivas algebraicamente; no se suman kVA indiscriminadamente ni se usa un único factor de demanda residencial para una central.

\section{Servicios propios propuestos}
El siguiente inventario es una \textbf{hipótesis de dimensionamiento}, con potencias eléctricas absorbidas. Para una bomba real cuya placa indique potencia mecánica se dividirá entre su eficiencia; no se duplicará esa corrección cuando el dato ya sea potencia eléctrica.
\small\input{generado/auxiliares}\normalsize
\small\input{generado/modos_auxiliares}\normalsize
En generación se consideran ocho luminarias interiores, cuatro exteriores, control, cargador, dos ventiladores, una bomba de drenaje y un gobernador activos: $P=\PAuxNormal$ kW y $Q=\QAuxNormal$ kvar. No se descuenta intermitencia de bombas sin un ciclo comprobado. La simultaneidad es una hipótesis operativa explícita, no un factor de demanda residencial. Mantenimiento suma $P=\PAuxMax$ kW y $S=\SAuxMax$ kVA con tomas e izaje, y se estudia con grupo parado e importación desde red.

Se conserva una bomba de reserva. Contar solo una exige enclavamiento y comprobar que la inundación de diseño no demanda ambas. La regla CNE-U 050-106(3) permite tratar cargas mutuamente excluyentes según la mayor, cuando esa exclusión está demostrada; 050-106(5) exige atender la máxima carga que pueda servirse para equipos cíclicos. No se reduce ampacidad simplemente por el consumo energético promedio.

El TSA candidato de 15 kVA alimenta auxiliares desde barra BT. Se suponen $P_{0,TSA}=0.08$ kW y $P_{cu,TSA}=0.25$ kW a placa. En generación sus pérdidas son \PerdTsa\ kW y se añaden al balance, sin confundirlas con carga de usuario. Sus potencias reactivas de magnetización y dispersión faltan en el modelo y deben incorporarse al flujo AC. La garantía en altura, corriente de arranque y caída transitoria de tensión pueden exigir otra potencia. Potencia continua suficiente no demuestra capacidad de arrancar el mayor motor.

\subsection{Arranque y transición entre modos}
Con red disponible y generador separado por 52G, el transformador principal puede alimentar la barra y el TSA si su esquema permite energización inversa. El control verifica agua, auxiliares, tensión, fase, velocidad y excitación antes de sincronizar. Se necesita enclavamiento que impida cierre fuera de sincronismo y recierre de red sobre una isla alimentada.

Al perder red, la CC debe mantener relés y disparos, pero una batería de control no demuestra capacidad para operar toda la unidad hidráulica del gobernador. Debe comprobarse acumulador hidráulico, retorno a posición segura o respaldo requerido para cerrar sin sobrepresión. El arranque autónomo queda sin demostrar: exige una fuente de servicios propios y una secuencia separada, además de control del sistema aislado.

\section{Energía anual y potencia firme}
La producción debe evaluarse con rendimientos variables y disponibilidad:
\begin{equation}
E_{net}=\sum_i\left[\rho g Q_{turb,i}H_{n,i}\eta_t\eta_m\eta_g/1000
-P_{aux,i}-P_{T,i}-P_{L,i}\right]\Delta t_i.
\end{equation}
La expresión se aplica en intervalos de operación; se modela por separado importación de auxiliares durante parada. Las horas fuera de servicio, vertimientos, límite del generador, restricciones de red y mínimos de turbina deben incorporarse. Si el despacho permite absorción de reactivos cambia la corriente y las pérdidas.

Como mera sensibilidad sobre el objetivo, $150\times8760\times FC$ da 657, 920 y 1117 MWh/año para $FC=0.50$, 0.70 y 0.85, redondeados. Ninguno es pronóstico del sitio. La energía firme, el factor de planta y la disponibilidad se mantienen sin valor validado hasta cerrar el recurso y operación.

\section{Campaña y datos de cierre}
\begin{enumerate}
\item Topografía georreferenciada y niveles de operación, incluyendo avenidas y restitución.
\item Series y aforos; inventario de derechos, usos y retornos; estudio ecológico según N05-N07.
\item Granulometría, concentración y mineralogía de sedimentos; temperatura del agua.
\item Registro horario de demanda o condiciones formales de exportación a red.
\item Inventario de auxiliares y garantías de desempeño del grupo para el rango de caudales.
\end{enumerate}
